\documentclass[10pt,a4paper]{article}
\usepackage[utf8]{inputenc}
\usepackage[T1]{fontenc}
\usepackage[spanish]{babel}
\usepackage{lmodern,amsmath,amssymb,booktabs,tabularx,array}
\usepackage[margin=1.85cm]{geometry}
\usepackage[hidelinks]{hyperref}
\setlength{\parindent}{0pt}
\setlength{\parskip}{4pt}
\newcommand{\R}{\mathbb{R}}
\begin{document}
\begin{center}
{\Large Reparación, aceptación y decisión numérica}\\[3pt]
{\large Seguimiento prospectivo independiente del paper OMS}\\[3pt]
14 de octubre de 2026 --- documento de investigación
\end{center}
\textbf{Resultado.} Verificar y reparar recuperó entregas que fallaban en Base,
pero no se demuestra una ventaja propia de contratos frente a alternativas
existentes. El estudio queda separado del artículo v1.3.5.
\section*{Diseño y aceptación}
Se fijaron 80 tratamientos antes de resolver: cinco modelos, semillas 11 y 23,
ocho alternativas y orden aleatorio registrado. Kondili H14/H16 son horizontes
nuevos de la misma fuente independiente; los tres MIPLIB son los del piloto,
con semillas nuevas. No son nuevas familias ni una muestra representativa.
La calibración previa, semilla 17, se excluye de la confirmación.

Gurobi 13.0.2, cuatro hilos, presolve y escalado interno automáticos, gap 1\%,
MILP hasta 15 s y referencia de ciclo de 20 s. Los LP condicionales usan
HiGHS 1.15.1 (hasta 2 s) y SoPlex 8.1.0/GMP (hasta 3 s restantes).
Las políticas escaladas usan exactamente el kernel diádico existente.

Se conservan los presupuestos: STN $10^{-5}$ kg (10 mg) en material,
$10^{-5}$ en lógica y $10^{-3}$ en unidades monetarias; MIPLIB
$10^{-4}\max(1,|b_r|)$ en unidades algebraicas experimentales.
Son objetivos numéricos del estudio. Una entrega debe satisfacer todas las
filas y cotas dentro de sus presupuestos y dominios enteros exactamente.
Los contratos MIPLIB cubren el punto crudo; el redondeado se comprueba después.
\begin{center}\small
\begin{tabular}{@{}lrrrrr@{}}\toprule
Tratamiento & Fuente & Entrega & LP & Abst. & Suma s\\\midrule
Default-check & 7/10 & 7/10 & 0 & 0 & 40.381\\
Tight-check & 10/10 & 10/10 & 0 & 0 & 48.326\\
NumericFocus-check & 6/10 & 6/10 & 0 & 0 & 32.690\\
IntegralityFocus-check & 10/10 & 10/10 & 0 & 0 & 68.679\\
Base-check-repair & 7/10 & 10/10 & 3 & 0 & 39.193\\
L2-check-repair & 7/10 & 10/10 & 3 & 0 & 52.895\\
Budget-check-repair & 6/10 & 8/10 & 2 & 2 & 32.765\\
Diagnostic-check-repair & 7/10 & 10/10 & 3 & 0 & 41.361\\
\bottomrule\end{tabular}\end{center}
Fuente: punto fuente tras redondear; Entrega: asignación verificada; LP:
reparaciones; Abst.: incompatibilidad contractual. Denominador diez por
tratamiento, incluyendo abstenciones. No son tasas poblacionales.

Hubo 78 llamadas MILP y 71
entregas verificadas. Los 11 LP flotantes condicionales recuperaron
la entrega, con costo de 0.0407--0.1760 s. No se llamó al LP exacto en
confirmación; dos puntos de HiGHS también resultaron factibles exactos.

La suma incluye selección, exportación, importación y cotejo, MILP,
verificación y reparación. La persistencia exterior se mide aparte; los
archivos auxiliares del LP están incluidos en reparación. No hubo errores
ni excesos de 20 s. \textbf{Desviación documentada:} la construcción común
inicial no se cronometró, pese a una frase del protocolo. No se incluye en
estas sumas. Dos semillas no permiten inferir velocidad general.
\newpage
\section*{Precisión, costo y calidad del objetivo}
En \texttt{control30-3-2-3}, Base-repair recupera precisión desde objetivos
crudos no entregables de $-39.917360$ y $-23.049913$, pero entrega costo
$26.248128$. Los controles estrictos entregan aproximadamente $1.27$ y
usan unos 15 s. La reparación rápida de peor objetivo no equivale a
optimización rápida de igual calidad. No se transfiere el gap de la fuente.
\begin{center}\small
\begin{tabular}{@{}rlrr@{}}\toprule
Semilla & Tratamiento & Costo entregado & Ciclo s\\\midrule
11 & Tight-check & 1.265274 & 15.154\\
11 & IntegralityFocus-check & 1.265272 & 15.114\\
11 & Base-check-repair & 26.248128 & 0.218\\
11 & Budget-check-repair & 24.661027 & 0.290\\
11 & L2-check-repair & 33.373658 & 0.295\\
23 & Tight-check & 1.265275 & 15.158\\
23 & IntegralityFocus-check & 1.315832 & 15.140\\
23 & Base-check-repair & 26.248128 & 0.242\\
23 & Budget-check-repair & 24.661027 & 0.230\\
23 & L2-check-repair & 33.373658 & 0.207\\
\bottomrule\end{tabular}\end{center}
En Kondili todas las alternativas pasan sin reparación. En cinca, Base,
Tight e IntegralityFocus pasan y NumericFocus falla en ambas semillas.
L2 necesita reparar una entrega; Budget se abstiene por 96 intervalos
incompatibles. En \texttt{ta1-UUM}, Base necesita reparación en semilla 23,
con cambio de objetivo de 0.282 unidades. La recuperación es del LP,
no del contrato. Los 80 registros y tiempos por etapa quedan conservados.
\section*{Contradicciones independientes de cotas}
Un punto binary64 reparado de \texttt{ta1-UUM} satisface exactamente todas
las filas, cotas y dominios del modelo canónico, comprobados con aritmética
racional. Su objetivo constituye una cota superior factible del mismo
modelo. Las llamadas Tight terminan con etiqueta óptima y cotas inferiores
que la contradicen:
\begin{center}\small
\begin{tabular}{@{}rrrr@{}}\toprule
Semilla & Cota reportada & Testigo exacto & Diferencia\\\midrule
11 & 7813164.58 & 7695229.19 & 117935.39\\
23 & 7842648.07 & 7695229.19 & 147418.88\\
\bottomrule\end{tabular}\end{center}
El testigo también mejora sus objetivos en más de 1\%. La importación y
el escalado se cotejan contra los mismos datos originales: no se mezclan
modelos. La guarda numérica del checker, $\max(10^{-6},8\,\mathrm{ulp}(B))$,
es muy inferior a las diferencias observadas. No se certifica el óptimo
global ni se consideran válidas las otras cotas por no hallar contradicción.
Tampoco se atribuye una causa al presolve o escalado interno sin controles
adicionales. Precisión primal, calidad y optimalidad se conservan separadas.
\newpage
\section*{Mecanismo de reparación y garantías de aceptación}
Los tres MIPLIB tienen máximos coeficientes enteros de $1.058\,10^7$,
$10^{17}$ y $1.109\,10^7$, frente a máximos continuos de 1, 1.4 y 1.
Fijar enteros retira los coeficientes grandes de la matriz del LP. Esta
observación orienta la reparación, pero no garantiza buena condición,
factibilidad del patrón o éxito del LP.

El ciclo selecciona y exporta, resuelve una vez y comprueba el punto
redondeado. Si falla, fija los enteros más próximos, reoptimiza el LP y
comprueba todas las restricciones originales. Si aún falla la entrega,
intenta un LP racional y comprueba por separado el testigo y su conversión
a binary64. Si nada pasa, se abstiene.

\textbf{Propiedad de aceptación.} Defina $V_\delta(y)$ como la verificación
de todas las filas y cotas bajo sus presupuestos originales y dominios
enteros exactos. Cada instrucción que devuelve un punto se ejecuta solamente
después de comprobar $V_\delta(y)$. Por tanto, todo punto devuelto satisface
$V_\delta$. Es una garantía del filtro, independiente de las etiquetas
del solver. No garantiza encontrar un punto ni certificar su optimalidad.
El número de intentos está acotado por un MILP y dos LP; la terminación
depende de que los backends terminen o sean interrumpidos. Los 20 s son
una referencia operativa, no una garantía de tiempo real.

Se sustituyen los enteros en el RHS con aritmética racional sobre los
datos binary64 originales. El LP racional se escribe con decimales completos,
sin un nuevo redondeo del RHS. Una contradicción en una fila constante
prueba ausencia de continuación exacta de ese patrón, no inviabilidad
global MILP. Un testigo exacto es primal; no certifica optimalidad MILP.

El desarrollo separado obtuvo un testigo exacto para control, con costo
26.248128, entrega flotante también verificada y tiempo
0.1091 s. Es desarrollo, excluido de la tabla prospectiva.

\section*{Obstrucción de representación}
Tras fijar enteros, considere $a x=b$, $a\ne0$, y $q=b/a$. Sean
$f^-<q<f^+$ los dos vecinos binary64 adyacentes, ambos finitos. Entonces
\[
\min_{f\in\mathbb F_{64}}|af-b|
=\min\{|af^--b|,|af^+-b|\}.
\]
En cada lado de $q$, $|a|\,|f-q|$ crece al alejarse; ningún otro valor
flotante tiene menor residuo. Si ese mínimo excede $\delta$, la fila no
tiene entrega binary64 que cumpla su presupuesto para esos enteros,
aunque admita un testigo racional exacto. Las cotas de variables no pueden
eliminar esa imposibilidad. El diagnóstico omite casos sin vecinos finitos.

Ejemplo de comprobación: $3x=1$, $0\le x\le1$. El testigo racional es
$1/3$ y el mínimo residuo binary64 es $2^{-54}$. Con $\delta=2^{-55}$
no hay entrega; con $\delta=2^{-54}$ sí. Cinco tests incluyen esa frontera,
contradicciones enteras, dominios y decimales exactos. No hubo obstrucciones
unarias en los casos reales. Es cuantización elemental, sin novedad
matemática reclamada ni beneficio práctico general demostrado.
\newpage
\section*{Diagnóstico previo y exportación}
La regla congelada eligió Budget para STN y Base con reparación reservada
para MIPLIB. Logra 10/10, como Base-repair, con mayor trabajo de selección.
En 48 puntos fuente que pasan, la cota nominal $\mu$ supera uno. La falta
de garantía suficiente no predice un fallo. Solo 12 puntos tienen
$\mu\le1$; el checker comprueba las hipótesis exportadas e integrales
y no encuentra contradicciones cuando son satisfechas.

El tratamiento del error de exportación, ya implementado en el piloto,
usa el soporte firmado de la caja original:
\[
R^+=\sum_j\max\{(\alpha a_j-\hat a_j)l_j,
(\alpha a_j-\hat a_j)u_j\}-(\alpha b-\hat b).
\]
Si $x\in[l,u]$ y $\hat a^\top x-\hat b\le\varepsilon$, descomponer
el residuo da
\[
a^\top x-b\le(\varepsilon+R^+)/\alpha.
\]
Si cada entero cambia a lo sumo $\gamma$ al redondear, la violación final
está acotada por
\[
\max\left\{0,\frac{\varepsilon+R^+}{\alpha}
+\gamma\sum_{j\in I}|a_j|\right\}.
\]
Para igualdades se consideran ambos signos. Términos con error cero no
necesitan cotas finitas. Aquí las exportaciones diádicas son exactamente
proporcionales y $R^+=0$. Se distinguen límite nominal, hipótesis observadas
y residuo efectivamente verificado, sin atribuir capacidad predictiva.
\section*{Antecedentes y decisión para OMS}
La reparación exacta y los certificados tienen precedentes explícitos:
Eifler y Gleixner (2021) y Cheung, Gleixner y Steffy (2016/2017).
SoPlex implementa LP racional e IntegralityFocus es un control existente
pertinente. Verificar o reparar no se reclama como idea novedosa.
La posible investigación es integrar requisitos, representación, aceptación,
abstención y calidad con trazabilidad. La muestra actual no demuestra una
ventaja propia del contrato frente a esos controles.

\textbf{Decisión: conservar el envío v1.3.5.} Quedan un ciclo realmente
ejecutado, reparaciones, certificados reproducibles y una separación
comprobada entre precisión y optimalidad. No se añade una contribución
ni una tabla al paper con estos resultados, ni se buscan casos favorables.
\section*{Reproducción y fuentes}
El checker independiente no importa el runner ni los solvers. Recomprueba
78 fuentes, 71 entregas, 11 candidatos LP, 9975 factores y 96 intervalos
incompatibles; coteja también los tres MPS originales y el desarrollo. Python/NumPy y
pytest bastan para replay y tests. Los programas, protocolos, enmienda
pre-ejecución, modelos, fuentes originales, logs, tiempos, puntos y hashes
se conservan en el artefacto. Los certificados se refieren al modelo
canónico binary64, no a decimales MPS racionales ideales o estándares físicos.

\small
\href{https://arxiv.org/abs/2101.09141}{Eifler y Gleixner: exact rational MIP (2021).}\\
\href{https://arxiv.org/abs/1611.08832}{Cheung, Gleixner y Steffy: verifying integer programming results.}\\
\href{https://soplex.zib.de/doc/html/EXACT.php}{SoPlex: exact LP and references to Gleixner--Steffy.}\\
\href{https://docs.gurobi.com/projects/optimizer/en/current/reference/parameters.html}{Gurobi: IntegralityFocus and numerical controls.}\\
\href{https://github.com/jckantor/STN-Scheduler}{STN-Scheduler: Kondili source, pinned MIT generator.}\\
\href{https://miplib.zib.de/tag_numerics.html}{MIPLIB 2017: Numerics collection and original MPS sources.}
\end{document}
